% Part: computability
% Chapter: recursive-functions
% Section: halting-problem

\documentclass[../../../include/open-logic-section]{subfiles}

\begin{document}

\olfileid{cmp}{rec}{hlt}
\olsection{थांबण्याची समस्या}

सर्वसाधारणपणे \emph{थांबण्याची समस्या} अशी आहे: संगणनक्षम फलनाचे
विनिर्देशन~$e$ (उदा., कार्यक्रम) आणि संख्या~$n$ दिली असता, आदान~$n$
वरील त्या फलनाचे संगणन थांबते का, म्हणजे काही परिणाम देते का, हे
ठरवणे. ही समस्या स्वतः संगणनक्षम फलनाने सोडवता येत नाही, हे ॲलन
ट्यूरिंग यांनी सिद्ध केले. म्हणजे, पुढील फलन संगणनक्षम नाही:
\[
h(e, n) =
\begin{cases}
1 & \text{जर संगणन $e$ आदान $n$ वर थांबत असेल तर}\\
0 & \text{अन्यथा,}
\end{cases}
\]

आंशिक पुनरावर्ती फलनांच्या संदर्भात कार्यक्रमाच्या विनिर्देशनाची भूमिका
क्लिनीच्या प्रमाण रूप प्रमेयातील निर्देशांक~$e$ बजावू शकतो. $f$ हे
आंशिक पुनरावर्ती फलन असेल, तर प्रमाण रूप प्रमेयातील समीकरण ज्या
कोणत्याही $e$~साठी खरे असेल, तो~$f$ चा निर्देशांक असतो. संख्या~$e$
दिली असता, प्रमाण रूप प्रमेय सांगते की
\[
\cfind{e}(x) \simeq U(\mu s \; T(e, x, s))
\]
हे आंशिक पुनरावर्ती असते; आणि प्रत्येक आंशिक पुनरावर्ती $f\colon \Nat
\to \Nat$ साठी असा $e \in \Nat$ असतो की सर्व~$x \in \Nat$ साठी
$\cfind{e}(x) \simeq f(x)$. खरे तर अशा प्रत्येक $f$ साठी एकच नव्हे,
तर अनंत निर्देशांक~$e$ असतात. \emph{थांबण्याचे फलन}~$h$ पुढीलप्रमाणे
परिभाषित करतात:
\[
h(e, x) =
\begin{cases}
1 & \text{जर $\cfind{e}(x) \fdefined$ असेल तर}\\
0 & \text{अन्यथा.}
\end{cases}
\]
$\cfind{e}(x) \fundefined$ असेल तेव्हा $h(e, x) = 0$, हे लक्षात घ्या.
प्रमाण रूपातील सूत्राने प्रत्येक~$e$ साठी आंशिक पुनरावर्ती फलन मिळते;
म्हणून येथे निर्देशांक नसल्याची वेगळी अपवादस्थिती उद्भवत नाही.
%READERNOTE{OLCMP-012 — गोठवलेल्या इंग्रजीत काही e आंशिक पुनरावर्ती फलनाचा निर्देशांक नसू शकतो असे म्हटले आहे. पण आधीच दिलेल्या प्रमाण रूपातील T आणि U मुळे प्रत्येक नैसर्गिक e हा एका आंशिक पुनरावर्ती फलनाचा निर्देशांक आहे. मराठीत ही अपवादस्थिती काढून सिद्धतेतील तदनुरूप प्रकरणे दुरुस्त केली आहेत.}

\begin{thm}
\ollabel{thm:halting-problem}
थांबण्याचे फलन~$h$ आंशिक पुनरावर्ती नाही.
\end{thm}

\begin{proof}
$h$ आंशिक पुनरावर्ती आहे असे समजल्यास, आपण पुढील फलन परिभाषित
करू शकू:
\[
d(y) =
\begin{cases}
1 & \text{जर $h(y, y) = 0$ असेल तर}\\
\umin{x}{x \neq x} & \text{अन्यथा.}
\end{cases}
\]
कोणतीही संख्या $x$ ही $x \neq x$ पूर्ण करत नसल्यामुळे $\umin{x}{x \neq
x}$ अस्तित्वात नाही. म्हणून $d(y) \fundefined$ तेव्हा आणि केवळ तेव्हाच,
जेव्हा $h(y,y) \neq 0$. या व्याख्येवरून पुढील गोष्टी मिळतात:
\begin{enumerate}
\item $d(y) \fdefined$ तेव्हा आणि केवळ तेव्हाच, जेव्हा $\cfind{y}(y)
  \fundefined$; प्रत्येक $y$ हा काही आंशिक पुनरावर्ती फलनाचा निर्देशांक
  असल्यामुळे वेगळी अपवादस्थिती नाही.
\item $d(y) \fundefined$ तेव्हा आणि केवळ तेव्हाच, जेव्हा $\cfind{y}(y)
  \fdefined$.
\end{enumerate}
$h$ आंशिक पुनरावर्ती असेल, तर $d$ ही आंशिक पुनरावर्ती असेल. त्यामुळे
क्लिनीच्या प्रमाण रूप प्रमेयानुसार तिला निर्देशांक~$e_d$ असेल.
$h(e_d, e_d)$ चे मूल्य पाहू. $0$ आणि~$1$ असे दोन पर्याय आहेत.
\begin{enumerate}
\item $h(e_d, e_d) = 1$ असेल, तर $\cfind{e_d}(e_d) \fdefined$.
  पण $\cfind{e_d} \simeq d$ आणि $d(e_d)$ हे $h(e_d, e_d) = 0$
  तेव्हा आणि केवळ तेव्हाच परिभाषित असते. म्हणून $h(e_d, e_d) \neq 1$.
\item $h(e_d, e_d) = 0$ असेल, तर $e_d$ हा त्या फलनाचा निर्देशांक असल्याने
  $\cfind{e_d}(e_d) \fundefined$. पण पुन्हा $\cfind{e_d} \simeq d$
  आणि $d(e_d)$ तेव्हा आणि केवळ तेव्हाच अपरिभाषित असते, जेव्हा
  $\cfind{e_d}(e_d) \fdefined$. त्यामुळे या प्रकरणातही विरोधाभास मिळतो.
\end{enumerate}
यावरून $e_d$ हा आंशिक पुनरावर्ती फलनाचा निर्देशांक असताना वरील
दोन्ही पर्याय अशक्य ठरतात. पण $h$ आंशिक पुनरावर्ती असेल, तर $d$ ही
आंशिक पुनरावर्ती असेल आणि तिचा निर्देशांक म्हणून $e_d$ निवडणे योग्य
असेल. म्हणून $h$ आंशिक पुनरावर्ती असू शकत नाही.
\end{proof}

\end{document}

